% ═══════════════════════════════════════════════════════════════
% LERNBLATT (Best-Practice-Design)
% Spanisch Klasse 8 - Repaso Unidad 5 + Números 100-1000
% ═══════════════════════════════════════════════════════════════
% DESIGN-PRINZIPIEN:
% - Sans-Serif (Helvetica): Fließtext
% - Serif (Palatino): Spanische Vokabeln (visuelle Distinktion)
% - Farbe NUR didaktisch begründet (Farbvokabeln)
% - Minimale Rahmen (kein "visual noise")
% ═══════════════════════════════════════════════════════════════

\documentclass[11pt, a4paper]{article}

\usepackage[utf8]{inputenc}
\usepackage[T1]{fontenc}
\usepackage[ngerman]{babel}

% --- SCHRIFTEN ---
\usepackage{helvet}
\renewcommand{\familydefault}{\sfdefault}
\usepackage{palatino}
\newcommand{\seriftext}[1]{{\fontfamily{ppl}\selectfont #1}}

\usepackage{microtype}
\usepackage[margin=1.8cm, top=2cm, bottom=1.5cm]{geometry}
\usepackage{enumitem}
\usepackage{tabularx}
\usepackage{booktabs}
\usepackage{array}
\usepackage{multicol}
\usepackage{tikz}

\usepackage{fancyhdr}
\usepackage{lastpage}

% ═══════════════════════════════════════════════════════════════
% FARBEN
% ═══════════════════════════════════════════════════════════════

\usepackage{xcolor}

% Farbvokabeln (didaktisch notwendig!)
\definecolor{colorblanco}{HTML}{FFFFFF}
\definecolor{colornegro}{HTML}{222222}
\definecolor{colorrojo}{HTML}{C62828}
\definecolor{colorazul}{HTML}{1565C0}
\definecolor{colorverde}{HTML}{2E7D32}
\definecolor{coloramarillo}{HTML}{F9A825}
\definecolor{colormarron}{HTML}{6D4C41}
\definecolor{colorgris}{HTML}{757575}

% Strukturfarben
\definecolor{hellgrau}{HTML}{F5F5F5}
\definecolor{rahmen}{HTML}{CCCCCC}
\definecolor{akzent}{HTML}{666666}

% ═══════════════════════════════════════════════════════════════
% BOXEN
% ═══════════════════════════════════════════════════════════════

\usepackage{tcolorbox}
\tcbuselibrary{skins}

\newtcolorbox{lernzielbox}{
    colback=hellgrau, colframe=hellgrau, boxrule=0pt, arc=0pt,
    left=10pt, right=10pt, top=8pt, bottom=8pt,
    before skip=6pt, after skip=8pt
}

\newtcolorbox{grammatikbox}[1][]{
    colback=white, colframe=white, boxrule=0pt, arc=0pt,
    left=10pt, right=6pt, top=6pt, bottom=6pt,
    borderline west={3pt}{0pt}{akzent},
    before skip=4pt, after skip=4pt
}

\newtcolorbox{merkbox}[1][]{
    colback=hellgrau, colframe=hellgrau, boxrule=0pt, arc=0pt,
    left=10pt, right=10pt, top=8pt, bottom=8pt,
    borderline north={2pt}{0pt}{colornegro},
    before skip=6pt, after skip=6pt
}

% ═══════════════════════════════════════════════════════════════
% BEFEHLE
% ═══════════════════════════════════════════════════════════════

\newcommand{\es}[1]{\seriftext{\textbf{#1}}}
\newcommand{\de}[1]{\textit{\textcolor{akzent}{#1}}}
\newcommand{\gram}[1]{\textsc{#1}}
\newcommand{\abmark}[1]{{\footnotesize\textcolor{akzent}{$\rightarrow$ AB #1}}}

% Farbmuster (quadratisch, 8pt)
\newcommand{\farbmuster}[1]{%
    \tikz[baseline=-0.5ex]\fill[#1] (0,0) rectangle (8pt,8pt);%
}

\setlength{\parindent}{0pt}
\setlength{\parskip}{5pt}

% ═══════════════════════════════════════════════════════════════
% KOPF- UND FUSSZEILE
% ═══════════════════════════════════════════════════════════════

\pagestyle{fancy}
\fancyhf{}
\renewcommand{\headrulewidth}{0.5pt}
\renewcommand{\footrulewidth}{0pt}
\lhead{\footnotesize\textsc{Spanisch} · Klasse 8}
\rhead{\footnotesize De compras}
\cfoot{\footnotesize\thepage\,/\,\pageref{LastPage}}

% ═══════════════════════════════════════════════════════════════
% DOKUMENT
% ═══════════════════════════════════════════════════════════════

\begin{document}

% --- TITEL ---
\begin{center}
    {\LARGE\bfseries De compras}\\[3pt]
    {\small Einkaufen · Kleidung · Farben · Zahlen bis 1000}
\end{center}

\vspace{4pt}

% --- EINLEITUNG ---
Du bist in Madrid und möchtest dir ein neues Outfit kaufen.
Der Verkäufer fragt: \es{¿En qué puedo ayudarte?} -- \de{Wie kann ich dir helfen?}

{\small\textit{Rate mal: Was kostet ein Wintermantel in Madrid -- 100\,€, 300\,€ oder 500\,€?
Die Antwort findest du auf Seite 2!}}

\vspace{2pt}

% --- LERNZIELE ---
\begin{lernzielbox}
\textbf{Das lernst du:}\quad
Kleidungsstücke benennen · Farben mit korrekter Endung · Modalverben konjugieren · Zahlen 100--1000 · Einkaufsdialog führen
\end{lernzielbox}

% ═══════════════════════════════════════════════════════════════
% ZWEI SPALTEN
% ═══════════════════════════════════════════════════════════════

\begin{multicols}{2}

% ---------------------------------------------------------------
% 1. KLEIDUNG
% ---------------------------------------------------------------
\textbf{\large 1. La ropa} \de{-- Kleidung}

\vspace{4pt}

\begin{tabular}{@{}ll@{}}
\es{la camiseta} & \de{das T-Shirt} \\[0.4em]
\es{el jersey} & \de{der Pullover} \\[0.4em]
\es{los pantalones} & \de{die Hose} \\[0.4em]
\es{el vestido} & \de{das Kleid} \\[0.4em]
\es{los zapatos} & \de{die Schuhe} \\[0.4em]
\es{la falda} & \de{der Rock} \\[0.4em]
\es{la chaqueta} & \de{die Jacke} \\[0.4em]
\es{el abrigo} & \de{der Mantel} \\
\end{tabular}

\vspace{4pt}
{\footnotesize\textit{Tipp:} \es{Los pantalones} steht immer im Plural -- genau wie im Englischen \textit{pants}!}

{\footnotesize\textit{Fun Fact:} \es{Jersey} kommt von der britischen Insel Jersey, wo die berühmten Wollpullover erfunden wurden.}

\abmark{1}

\vspace{8pt}

% ---------------------------------------------------------------
% 2. FARBEN (mit Farbmustern)
% ---------------------------------------------------------------
\textbf{\large 2. Los colores} \de{-- Farben}

\vspace{4pt}

\begin{tabular}{@{}cll@{}}
\farbmuster{colorblanco} & \es{blanco/a} & \de{weiß} \\[0.4em]
\farbmuster{colornegro} & \es{negro/a} & \de{schwarz} \\[0.4em]
\farbmuster{colorrojo} & \es{rojo/a} & \de{rot} \\[0.4em]
\farbmuster{coloramarillo} & \es{amarillo/a} & \de{gelb} \\[0.4em]
\farbmuster{colorazul} & \es{azul} & \de{blau} \\[0.4em]
\farbmuster{colorverde} & \es{verde} & \de{grün} \\[0.4em]
\farbmuster{colormarron} & \es{marrón} & \de{braun} \\[0.4em]
\farbmuster{colorgris} & \es{gris} & \de{grau} \\
\end{tabular}

\vspace{6pt}

\begin{grammatikbox}
\textbf{\gram{Genus}-Anpassung}

Farben auf \es{-o/-a} ändern sich:\\
{\small\es{jersey negr\underline{o}} · \es{camiseta blanc\underline{a}}}

Farben auf \es{-e} / Konsonant bleiben:\\
{\small\es{vestido verde} · \es{falda verde}}
\end{grammatikbox}

\abmark{2}

\columnbreak

% ---------------------------------------------------------------
% 3. MODALVERBEN
% ---------------------------------------------------------------
\textbf{\large 3. Verbos modales} \de{-- Modalverben}

\vspace{4pt}

Mit Modalverben sagst du, was du \textbf{willst}, \textbf{kannst} oder \textbf{musst}:

\begin{center}
\es{Quiero comprar un jersey.}\\
\de{= Ich will einen Pullover kaufen.}
\end{center}

\vspace{2pt}

\begin{grammatikbox}
{\small
\begin{tabular}{@{}l@{\,}l@{\,}l@{\,}l@{}}
& \textbf{querer} & \textbf{poder} & \textbf{tener que} \\
& \de{wollen} & \de{können} & \de{müssen} \\[0.3em]
\midrule
yo & \es{qu\textcolor{colorrojo}{ie}ro} & \es{p\textcolor{colorrojo}{ue}do} & \es{tengo que} \\
tú & \es{qu\textcolor{colorrojo}{ie}res} & \es{p\textcolor{colorrojo}{ue}des} & \es{tienes que} \\
él/ella & \es{qu\textcolor{colorrojo}{ie}re} & \es{p\textcolor{colorrojo}{ue}de} & \es{tiene que} \\
nosotros & \es{queremos} & \es{podemos} & \es{tenemos que} \\
vosotros & \es{queréis} & \es{podéis} & \es{tenéis que} \\
ellos & \es{qu\textcolor{colorrojo}{ie}ren} & \es{p\textcolor{colorrojo}{ue}den} & \es{tienen que} \\
\end{tabular}}
\end{grammatikbox}

\vspace{2pt}
{\footnotesize\textit{Merke:} Stammwechsel \textcolor{colorrojo}{e→ie} und \textcolor{colorrojo}{o→ue} -- außer bei nosotros/vosotros!}

\abmark{3}

\vspace{8pt}

% ---------------------------------------------------------------
% 4. EINKAUFSDIALOG
% ---------------------------------------------------------------
\textbf{\large 4. De compras} \de{-- Einkaufen}

\vspace{4pt}

\textbf{Vendedor/a} \de{(Verkäufer/in):}\\
{\small
\es{¿En qué puedo ayudarte?}\\
\es{¿Qué color/talla quieres?}\\
\es{Cuesta ... euros.}}

\vspace{4pt}

\textbf{Cliente} \de{(Kunde/in):}\\
{\small
\es{Quiero comprar...}\\
\es{¿Cuánto cuesta?}\\
\es{Lo/La compro.} · \es{Es muy caro.}}

\vspace{6pt}

\textbf{Beispieldialog:}\\
{\footnotesize
\textbf{V:} \es{¡Hola! ¿En qué puedo ayudarte?}\\
\textbf{C:} \es{Quiero comprar un jersey \textcolor{colorazul}{azul}.}\\
\textbf{V:} \es{Cuesta 45 euros.}\\
\textbf{C:} \es{Bien, lo compro.}}

\abmark{6}

\end{multicols}

% ═══════════════════════════════════════════════════════════════
% SEITENUMBRUCH
% ═══════════════════════════════════════════════════════════════
\newpage

% ---------------------------------------------------------------
% 5. ZAHLEN 100-1000
% ---------------------------------------------------------------
\textbf{\large 5. Los números 100--1000} \de{-- Zahlen}

\vspace{4pt}

Preise in Spanien sind oft dreistellig. Ein guter Wintermantel kostet ca. \es{trescientos euros} -- das war die Lösung vom Rätsel!

\vspace{4pt}

\begin{merkbox}
\textbf{Die Hunderter}

\vspace{4pt}
\begin{tabular}{@{}rl@{\hspace{0.6cm}}rl@{\hspace{0.6cm}}rl@{\hspace{0.6cm}}rl@{}}
100 & \es{cien} & 300 & \es{trescientos} & 600 & \es{seiscientos} & 900 & \es{novecientos} \\[0.3em]
200 & \es{doscientos} & 400 & \es{cuatrocientos} & 700 & \es{setecientos} & 1000 & \es{mil} \\[0.3em]
& & 500 & \es{quinientos} & 800 & \es{ochocientos} & & \\
\end{tabular}

\vspace{6pt}

% Zahlenstrahl mit roten Punkten AUF der Linie
\begin{center}
\begin{tikzpicture}[scale=0.9]
    \draw[thick, ->] (0,0) -- (11,0);
    \foreach \x/\num in {0/100, 1/200, 2/300, 3/400, 4/500, 5/600, 6/700, 7/800, 8/900, 9/1000} {
        \draw[thick] (\x,0.1) -- (\x,-0.1);
        \node[below] at (\x,-0.15) {\scriptsize\num};
    }
    % Unregelmäßige Zahlen: rote Punkte AUF dem Zahlenstrahl
    \fill[colorrojo] (4,0) circle (4pt);  % 500
    \fill[colorrojo] (6,0) circle (4pt);  % 700
    \fill[colorrojo] (8,0) circle (4pt);  % 900
\end{tikzpicture}

\vspace{2pt}
{\scriptsize\textcolor{colorrojo}{$\bullet$} = unregelmäßige Schreibweise}
\end{center}
\end{merkbox}

\vspace{4pt}

\begin{multicols}{2}

\begin{grammatikbox}
\textbf{Drei Regeln}

\textbf{1.} 100 allein = \es{cien}\\
\phantom{1.} 100+ = \es{ciento uno}, \es{ciento cincuenta}

\vspace{3pt}
\textbf{2.} Unregelmäßig {\small\textcolor{colorrojo}{$\bullet$}}:\\
\phantom{2.} \es{quinientos} · \es{setecientos} · \es{novecientos}\\
\phantom{2.} {\footnotesize Warum? Aus dem Latein: \textit{quingenti}, \textit{septingenti}!}

\vspace{3pt}
\textbf{3.} \gram{Genus}-Anpassung ab 200:\\
\phantom{3.} \es{doscientos euros} \de{(m.)}\\
\phantom{3.} \es{doscientas personas} \de{(f.)}
\end{grammatikbox}

\columnbreak

\textbf{Beispiele:}\\[2pt]
{\small
\begin{tabular}{@{}rl@{}}
350 & = \es{trescientos cincuenta} \\[0.3em]
725 & = \es{setecientos veinticinco} \\[0.3em]
999 & = \es{novecientos noventa y nueve} \\[0.3em]
156 & = \es{ciento cincuenta y seis} \\
\end{tabular}}

\vspace{8pt}

\textbf{Zusammengesetzt:}\\
{\small \es{ciento} + Zehner + \es{y} + Einer\\
z.B. 156 = \es{ciento cincuenta y seis}}

\vspace{6pt}

{\footnotesize
\textit{Basis:} Lerne zuerst die glatten Hunderter (100--500).\\
\textit{Fortgeschritten:} Übe zusammengesetzte wie 725, 999.}

\end{multicols}

\abmark{4 + 5}

\vspace{8pt}

% ═══════════════════════════════════════════════════════════════
% ZUSATZAUFGABEN
% ═══════════════════════════════════════════════════════════════

\begin{tcolorbox}[
    colback=white, colframe=rahmen, boxrule=0.5pt, arc=0pt,
    left=8pt, right=8pt, top=6pt, bottom=6pt,
    borderline north={1.5pt}{0pt}{akzent}
]
{\small\textbf{Zusatzaufgaben für Schnelle}}

\vspace{4pt}

\begin{multicols}{2}
\footnotesize

\textbf{Z1: Mein Outfit}\\
Beschreibe 5 Kleidungsstücke mit Farben und Preisen auf Spanisch.

\vspace{4pt}

\textbf{Z2: Partner-Dialog}\\
Schreibt einen Einkaufsdialog (8 Zeilen).

\columnbreak

\textbf{Z3: Rechnen}\\
Ergebnis als spanisches Wort:\\
a) 350 + 275 = \rule{2cm}{0.3pt}\\
b) 1000 $-$ 456 = \rule{2cm}{0.3pt}\\
c) 125 $\times$ 4 = \rule{2cm}{0.3pt}

\vspace{4pt}

\textbf{Z4: Kreativ}\\
Kaufe ein Geschenk. Schreibe Dialog.

\end{multicols}
\end{tcolorbox}

\vspace{8pt}

% ═══════════════════════════════════════════════════════════════
% WUSSTEST DU? (Kulturelles)
% ═══════════════════════════════════════════════════════════════

\begin{grammatikbox}
\textbf{¿Sabías que...?} \de{-- Wusstest du?}

\vspace{4pt}
\begin{multicols}{2}
\footnotesize

$\bullet$ In Spanien sagt man beim Bezahlen oft \es{¿Me cobra?} \de{(Können Sie kassieren?)} statt direkt nach dem Preis zu fragen.

\vspace{3pt}
$\bullet$ \es{Rebajas} \de{(Schlussverkauf)} gibt es zweimal im Jahr: Januar und Juli. Dann sind viele Geschäfte \es{de rebajas}!

\columnbreak

$\bullet$ Das Wort \es{euro} ist im Spanischen maskulin: \es{un euro}, \es{dos euros}. Aber: \es{cien euros} (nicht \es{ciento}!).

\vspace{3pt}
$\bullet$ \es{Lo/La compro} -- Das Pronomen muss zum Objekt passen: \es{el jersey} $\rightarrow$ \es{lo compro}, \es{la camiseta} $\rightarrow$ \es{la compro}.

\end{multicols}
\end{grammatikbox}

\vspace{6pt}

% Eselsbrücke für die unregelmäßigen Zahlen
\begin{center}
\footnotesize\textit{Merksatz für 500, 700, 900:} \textbf{``Quin-Set-Nov''} -- die drei mit lateinischen Wurzeln!
\end{center}

\end{document}
